\chapter{Additional Matrix Algebra}
\label{app:matrix}

This appendix collects material that is useful for reference but is not
needed before the main econometrics chapters.

\section{Addition and scalar multiplication}

Matrices can be added only when they have the same dimensions, and addition
is entry by entry.  If $A$ and $B$ are $m\times k$, then
\[
  (A+B)_{ij}=a_{ij}+b_{ij}.
\]
For a scalar $c$, scalar multiplication is also entry by entry:
$(cA)_{ij}=ca_{ij}$.  Consequently,
\[
  A+B=B+A,\qquad c(A+B)=cA+cB.
\]

\section{Several views of multiplication}

Let $A$ be $n\times k$, $B$ be $n\times p$, and $C=A'B$.  In addition to
the entry formula in Chapter~\ref{ch:matrix}, we can write
\[
  C_{ab}=A_{\cdot a}'B_{\cdot b}
  \quad\text{and}\quad
  A'B=\sum_{i=1}^{n}A_iB_i',
\]
where $A_{\cdot a}$ denotes column $a$ of $A$, while $A_i$ and $B_i$ are the
column-vector versions of row $i$.  Each term $A_iB_i'$ is $k\times p$.

The following code verifies the two representations.  It is pedagogically
useful, although in practice one should use R's optimized matrix product.

\begin{rcode}
A <- matrix(c(1, 2, 3, 4, 5, 6), nrow = 3, ncol = 2)
B <- matrix(c(2, 3, 4, 5, 6, 7), nrow = 3, ncol = 2)

# Direct product
C <- t(A) %*% B

# Entry-by-entry representation
C_entry <- matrix(0, ncol(A), ncol(B))
for (a in seq_len(ncol(A))) {
  for (b in seq_len(ncol(B))) {
    C_entry[a, b] <- sum(A[, a] * B[, b])
  }
}

# Sum of row outer products
C_outer <- matrix(0, ncol(A), ncol(B))
for (i in seq_len(nrow(A))) {
  ai <- matrix(A[i, ], ncol = 1)
  bi <- matrix(B[i, ], nrow = 1)
  C_outer <- C_outer + ai %*% bi
}

C
C_entry
C_outer
\end{rcode}

\section{Rank and inverse formulas}

Some useful rank facts are
\[
  \operatorname{rank}(AB)
  \leq \min\{\operatorname{rank}(A),\operatorname{rank}(B)\}
\]
and
\[
  \operatorname{rank}(A'A)=\operatorname{rank}(AA')
  =\operatorname{rank}(A).
\]
The second equality does \emph{not} imply that both $A'A$ and $AA'$ are
invertible when $A$ is rectangular.  If $A$ is $m\times k$ with $m>k$ and
full column rank, then $A'A$ is invertible but $AA'$ has rank $k<m$ and is
singular.  The roles reverse when $m<k$ and $A$ has full row rank.

For an invertible diagonal matrix
$D=\operatorname{diag}(d_1,\ldots,d_k)$,
\[
  D^{-1}=\operatorname{diag}(d_1^{-1},\ldots,d_k^{-1}).
\]
For
\[
  A=\begin{pmatrix}a&b\\c&d\end{pmatrix},
\]
the determinant is $ad-bc$, and, when $ad-bc\neq0$,
\[
  A^{-1}=\frac{1}{ad-bc}
  \begin{pmatrix}d&-b\\-c&a\end{pmatrix}.
\]

\section{Eigenvalues and spectral decomposition}

An eigenvalue--eigenvector pair $(\lambda,v)$ of a square matrix $A$ satisfies
$Av=\lambda v$ with $v\neq0$.  A general real square matrix need not have a
full set of real eigenvectors and need not be diagonalizable.  For example,
$\begin{psmallmatrix}1&1\\0&1\end{psmallmatrix}$ is not diagonalizable.

The clean result needed most often in econometrics applies to symmetric
matrices.

\begin{theorem}[Spectral theorem]
For every real symmetric $k\times k$ matrix $A$, there are an orthogonal
matrix $H$ and a real diagonal matrix $\Lambda$ such that
\[
  A=H\Lambda H',\qquad H'H=HH'=I_k.
\]
The diagonal entries of $\Lambda$ are the eigenvalues of $A$, and the columns
of $H$ are orthonormal eigenvectors.
\end{theorem}

The decomposition is not unique: eigenvectors can be reordered, their signs
can be changed, and within an eigenspace associated with a repeated
eigenvalue, any orthonormal basis may be used.

For a real symmetric matrix:
\begin{itemize}
  \item $A$ is p.s.d. if and only if every eigenvalue is nonnegative;
  \item $A$ is p.d. if and only if every eigenvalue is positive;
  \item $\operatorname{rank}(A)$ equals the number of nonzero eigenvalues;
  \item $A$ is invertible if and only if zero is not an eigenvalue, and then
        $A^{-1}=H\Lambda^{-1}H'$.
\end{itemize}
The rank statement is not valid for every nonsymmetric matrix.  A nilpotent
matrix such as $\begin{psmallmatrix}0&1\\0&0\end{psmallmatrix}$ has rank one
but only the eigenvalue zero.

For a rectangular matrix, the analogous tool is the singular-value
decomposition $A=U\Sigma V'$.  Its rank equals the number of nonzero singular
values.  This is why singular values, rather than ordinary eigenvalues, are
the safe general-purpose rank diagnostic.

\begin{rcode}
library(Matrix)
A <- matrix(c(2, 1, 1, 2), 2, 2)
eigen(A)
rankMatrix(A)
solve(A)

# A nonsymmetric matrix that is not diagonalizable
J <- matrix(c(1, 0, 1, 1), 2, 2)
eigen(J)
\end{rcode}

\section{Matrix derivatives}

We use the convention that the gradient of a scalar-valued function with
respect to a column vector is itself a column vector.  For conformable fixed
objects $x$, $y$, and $X$,
\[
  \nabla_\beta(x'\beta)=x,
  \qquad
  \nabla_\beta(\beta'\beta)=2\beta,
  \qquad
  \frac{\partial(X\beta)}{\partial\beta'}=X.
\]
The last object is a Jacobian, not a scalar gradient.

For the least-squares criterion
\[
  Q(\beta)=(Y-X\beta)'(Y-X\beta),
\]
expansion or the chain rule gives
\begin{align*}
  \nabla_\beta Q(\beta)
  &=-2X'(Y-X\beta)\\
  &=2X'(X\beta-Y).
\end{align*}
The minus sign is essential.  Setting the gradient to zero yields the normal
equations $X'X\beta=X'Y$, and if $X$ has full column rank,
$\widehat\beta=(X'X)^{-1}X'Y$.

\section{Further reading}

Useful extensions include determinants, traces, idempotent projection
matrices, generalized inverses, and singular-value decomposition.  A compact
reference is Petersen and Pedersen's \emph{Matrix Cookbook}.
